\section{第 14.3 节课程题目：自旋态、Pauli 代数与旋转}
\label{exercises:QM-C14-S03}
\exercisemeta
  {QM-C14-S03}
  {2026-10-29 至 2026-10-31}
  {Shankar, Second Edition；正文示范印刷 pp.~383--384 / PDF pp.~390--391；Exercises 14.3.2、14.3.4 在印刷 p.~384 / PDF p.~391，14.3.6 在印刷 p.~385 / PDF p.~392}
  {一处教材正文示范已讲解；三道选定习题已作答并核对}
  {笔记已确认}

本节无独立编号 Example；第一项整理教材正文中的旋转示范。习题 14.3.4 只选第一种证明方法，其余未选题不要求补做。以下补齐引用公式与约定，不收入课堂概念检查。

\exercisequestion{QM-C14-S03-E01}{教材正文示范：一次旋转得到任意方向的自旋态}
\textbf{来源：}印刷 pp.~383--384，式 (14.3.45)--(14.3.47)，无独立 Example 编号。

\textbf{对应知识点：}\relatedtopic{QM-C14-S03-T05}{有限旋转与半角}。

\textbf{示范任务。}初态为 $\ket{z,+}=(1,0)^{\mathsf T}$。目标方向为
$\hat{\mathbf n}=(\sin\theta\cos\varphi,\sin\theta\sin\varphi,\cos\theta)$。
先取 $0<\theta<\pi$，选择把 $+\hat{\mathbf z}$ 转到 $\hat{\mathbf n}$ 的旋转轴，用
\begin{equation}
  U_{\mathbf a}(\alpha)=\cos\frac\alpha2 I
    -\ii\sin\frac\alpha2\hat{\mathbf a}\cdot\boldsymbol\sigma
\end{equation}
求旋转后的态，并与教材相位约定
$\ket{\hat{\mathbf n},+}=(\cos(\theta/2)\ee^{-\ii\varphi/2},
\sin(\theta/2)\ee^{\ii\varphi/2})^{\mathsf T}$ 比较。

\begin{checkedsolution}
按右手定则选轴和角度：
\begin{equation}
  \hat{\mathbf a}=\frac{\hat{\mathbf z}\times\hat{\mathbf n}}
     {|\hat{\mathbf z}\times\hat{\mathbf n}|}
     =(-\sin\varphi,\cos\varphi,0),\qquad \alpha=\theta.
\end{equation}
由 Pauli 矩阵，
\begin{equation}
  \hat{\mathbf a}\cdot\boldsymbol\sigma
    =\begin{pmatrix}0&-\ii\ee^{-\ii\varphi}\\
      \ii\ee^{\ii\varphi}&0\end{pmatrix},\qquad
  U_{\mathbf a}(\theta)=\begin{pmatrix}
      \cos(\theta/2)&-\sin(\theta/2)\ee^{-\ii\varphi}\\
      \sin(\theta/2)\ee^{\ii\varphi}&\cos(\theta/2)
    \end{pmatrix}.
\end{equation}
作用于初态即取第一列：
\begin{equation}
  \boxed{U_{\mathbf a}(\theta)\ket{z,+}
    =\begin{pmatrix}\cos(\theta/2)\\\sin(\theta/2)\ee^{\ii\varphi}\end{pmatrix}
    =\ee^{\ii\varphi/2}\ket{\hat{\mathbf n},+}.}
\end{equation}
两者只差整体相位，物理状态相同。$\theta=0,\pi$ 时不能再用叉积归一化来定义轴，但上面的显式选轴与旋转矩阵仍适用；极点处方位角不再指定不同的目标方向。
\end{checkedsolution}

\exercisequestion{QM-C14-S03-E02}{教材习题 14.3.2：本征旋量与自旋期望值}
\textbf{对应知识点：}\relatedtopic{QM-C14-S03-T03}{任意方向的自旋态}。

\textbf{题目。}在 $S_z$ 本征基下取
\begin{equation}
  S_x=\frac\hbar2\begin{pmatrix}0&1\\1&0\end{pmatrix},\quad
  S_y=\frac\hbar2\begin{pmatrix}0&-\ii\\\ii&0\end{pmatrix},\quad
  S_z=\frac\hbar2\begin{pmatrix}1&0\\0&-1\end{pmatrix}.
\end{equation}
令 $\hat{\mathbf n}=(\sin\theta\cos\varphi,\sin\theta\sin\varphi,\cos\theta)$，$0\le\theta\le\pi$，并给出
\begin{equation}
  \chi_+=\begin{pmatrix}\cos(\theta/2)\ee^{-\ii\varphi/2}\\
      \sin(\theta/2)\ee^{\ii\varphi/2}\end{pmatrix},\qquad
  \chi_-=\begin{pmatrix}-\sin(\theta/2)\ee^{-\ii\varphi/2}\\
      \cos(\theta/2)\ee^{\ii\varphi/2}\end{pmatrix}.
\end{equation}
验证 $S_{\mathbf n}\chi_\pm=\pm\hbar\chi_\pm/2$，并分别计算三个分量的期望值，证明 $\expect{\mathbf S}_\pm=\pm\hbar\hat{\mathbf n}/2$。原题使用 $\hat{\mathbf n}\cdot\boldsymbol\sigma$ 的本征值 $\pm1$，与此完全等价。

\begin{checkedsolution}
\textbf{1. 本征态。}记 $c=\cos(\theta/2)$、$s=\sin(\theta/2)$。由
$\cos\theta=c^2-s^2$、$\sin\theta=2cs$，可得
\begin{equation}
  c\cos\theta+s\sin\theta=c,\qquad
  c\sin\theta-s\cos\theta=s.
\end{equation}
于是直接相乘给出
\begin{align}
  S_{\mathbf n}\chi_+
   &=\frac\hbar2\begin{pmatrix}
     \ee^{-\ii\varphi/2}(c\cos\theta+s\sin\theta)\\
     \ee^{\ii\varphi/2}(c\sin\theta-s\cos\theta)
     \end{pmatrix}=\frac\hbar2\chi_+,\\
  S_{\mathbf n}\chi_-
   &=\frac\hbar2\begin{pmatrix}
     \ee^{-\ii\varphi/2}(-s\cos\theta+c\sin\theta)\\
     \ee^{\ii\varphi/2}(-s\sin\theta-c\cos\theta)
     \end{pmatrix}=-\frac\hbar2\chi_-.
\end{align}
同时 $\chi_\pm^\dagger\chi_\pm=c^2+s^2=1$、$\chi_+^\dagger\chi_-=0$，故确为正交归一本征基。

\textbf{2. 期望值。}对任意归一化列向量 $\chi=(a,b)^{\mathsf T}$，直接矩阵相乘得到
\begin{equation}
  \expect{S_x}=\hbar\operatorname{Re}(a^*b),\qquad
  \expect{S_y}=\hbar\operatorname{Im}(a^*b),\qquad
  \expect{S_z}=\frac\hbar2(|a|^2-|b|^2).
\end{equation}
其中 $S_y$ 的组合是 $\hbar(-\ii a^*b+\ii b^*a)/2$，取伴随时相位须共轭。对 $\chi_+$ 有 $a^*b=cs\ee^{\ii\varphi}$，对 $\chi_-$ 有 $a^*b=-cs\ee^{\ii\varphi}$，所以
\begin{equation}
  \begin{array}{c|ccc}
    &\expect{S_x}&\expect{S_y}&\expect{S_z}\\ \hline
    \chi_+&\hbar cs\cos\varphi&\hbar cs\sin\varphi&\hbar(c^2-s^2)/2\\
    \chi_-&-\hbar cs\cos\varphi&-\hbar cs\sin\varphi&\hbar(s^2-c^2)/2
  \end{array}
\end{equation}
利用 $2cs=\sin\theta$，最终
\begin{equation}
  \boxed{\expect{\mathbf S}_\pm
    =\pm\frac\hbar2(\sin\theta\cos\varphi,\sin\theta\sin\varphi,\cos\theta).}
\end{equation}
这些是平均值，不表示三个分量同时具有确定测量值。
\end{checkedsolution}

\exercisequestion{QM-C14-S03-E03}{教材习题 14.3.4：由对易关系推导乘积恒等式}
\textbf{对应知识点：}\relatedtopic{QM-C14-S03-T04}{Pauli 矩阵代数}。

\textbf{题目。}只采用原题的第一种方法。设 $\mathbf A,\mathbf B$ 是普通实数矢量，已知
\begin{equation}
  [\sigma_i,\sigma_j]=2\ii\sum_k\epsilon_{ijk}\sigma_k,
  \qquad \{\sigma_i,\sigma_j\}=2\delta_{ij}I.
\end{equation}
先将 $\sigma_i\sigma_j$ 用对易子与反对易子表示，再证明
$(\mathbf A\cdot\boldsymbol\sigma)(\mathbf B\cdot\boldsymbol\sigma)
=(\mathbf A\cdot\mathbf B)I+\ii(\mathbf A\times\mathbf B)\cdot\boldsymbol\sigma$。

\begin{checkedsolution}
将对易子与反对易子相加：
\begin{equation}
  \sigma_i\sigma_j=\frac12\bigl(\{\sigma_i,\sigma_j\}+[\sigma_i,\sigma_j]\bigr)
    =\delta_{ij}I+\ii\sum_k\epsilon_{ijk}\sigma_k.
\end{equation}
于是
\begin{align}
  (\mathbf A\cdot\boldsymbol\sigma)(\mathbf B\cdot\boldsymbol\sigma)
    &=\sum_{i,j}A_iB_j\sigma_i\sigma_j\notag\\
    &=\sum_iA_iB_iI
       +\ii\sum_k\left(\sum_{i,j}\epsilon_{ijk}A_iB_j\right)\sigma_k.
\end{align}
第一项是 $(\mathbf A\cdot\mathbf B)I$。循环置换不改变 Levi-Civita 符号，故
\begin{equation}
  \epsilon_{ijk}=\epsilon_{kij},\qquad
  \sum_{i,j}\epsilon_{ijk}A_iB_j=(\mathbf A\times\mathbf B)_k.
\end{equation}
因此
\begin{equation}
  \boxed{(\mathbf A\cdot\boldsymbol\sigma)(\mathbf B\cdot\boldsymbol\sigma)
    =(\mathbf A\cdot\mathbf B)I
     +\ii(\mathbf A\times\mathbf B)\cdot\boldsymbol\sigma.}
\end{equation}
这里是循环置换，不是交换两个指标，所以没有额外负号；标量点积项前的单位矩阵也不可省略其含义。
\end{checkedsolution}

\exercisequestion{QM-C14-S03-E04}{教材习题 14.3.6：先绕固定 y 轴，再绕固定 z 轴}
\textbf{对应知识点：}\relatedtopic{QM-C14-S03-T05}{有限旋转与顺序}。

\textbf{题目。}初态为 $\ket{z,+}=(1,0)^{\mathsf T}$，目标方向为
$\hat{\mathbf n}=(\sin\theta\cos\varphi,\sin\theta\sin\varphi,\cos\theta)$。
采用右手定则的主动旋转，且两次均绕固定坐标轴。给定
\begin{equation}
  U_y(\theta)=\ee^{-\ii\theta\sigma_y/2},\quad
  U_z(\varphi)=\ee^{-\ii\varphi\sigma_z/2},\quad
  \sigma_y=\begin{pmatrix}0&-\ii\\\ii&0\end{pmatrix},\quad
  \sigma_z=\begin{pmatrix}1&0\\0&-1\end{pmatrix}.
\end{equation}
先从几何说明 $R_z(\varphi)R_y(\theta)\hat{\mathbf z}=\hat{\mathbf n}$ 及算符的作用顺序，再用 $\ee^{-\ii\alpha\sigma_i/2}=\cos(\alpha/2)I-\ii\sin(\alpha/2)\sigma_i$ 显式验证
\begin{equation}
  U_z(\varphi)U_y(\theta)\ket{z,+}
    =\begin{pmatrix}\cos(\theta/2)\ee^{-\ii\varphi/2}\\
      \sin(\theta/2)\ee^{\ii\varphi/2}\end{pmatrix}
    =\ket{\hat{\mathbf n},+}.
\end{equation}

\begin{checkedsolution}
\textbf{1. 几何与顺序。}先绕 $y$ 轴改变极角，再绕 $z$ 轴改变方位角：
\begin{equation}
  \begin{pmatrix}0\\0\\1\end{pmatrix}
    \xrightarrow{R_y(\theta)}
  \begin{pmatrix}\sin\theta\\0\\\cos\theta\end{pmatrix}
    \xrightarrow{R_z(\varphi)}
  \begin{pmatrix}\sin\theta\cos\varphi\\\sin\theta\sin\varphi\\\cos\theta\end{pmatrix}.
\end{equation}
\begin{center}
\begin{tikzpicture}[>=Stealth,scale=0.9,font=\small]
  \begin{scope}
    \draw[->] (0,0)--(2.55,0) node[right] {$x$};
    \draw[->] (0,0)--(0,2.5) node[above] {$z$};
    \draw[->,very thick,noteBlue] (0,0)--(0,2.1);
    \draw[->,very thick,accentRed] (0,0)--(40:2.1)
      node[above right] {$\mathbf n_1$};
    \draw[->] (0,0.72) arc[start angle=90,end angle=40,radius=0.72];
    \node at (0.39,0.9) {$\theta$};
    \node[align=center] at (1.05,-0.6) {先绕 $y$ 轴\\$xz$ 平面};
  \end{scope}
  \begin{scope}[xshift=6cm]
    \draw[->] (0,0)--(2.55,0) node[right] {$x$};
    \draw[->] (0,0)--(0,2.5) node[above] {$y$};
    \draw[->,very thick,noteBlue] (0,0)--(2.1,0);
    \draw[->,very thick,accentRed] (0,0)--(40:2.1)
      node[above right] {$\mathbf n_{xy}$};
    \draw[->] (0.72,0) arc[start angle=0,end angle=40,radius=0.72];
    \node at (0.98,0.35) {$\varphi$};
    \node[align=center] at (1.05,-0.6) {再绕 $z$ 轴\\仅画 $xy$ 投影};
  \end{scope}
\end{tikzpicture}

{\small 原创方向示意，角度与长度不按同一比例；第二步保持 $z$ 分量 $\cos\theta$ 不变。}
\end{center}
列向量从右向左受作用，所以先发生的 $U_y$ 写在右边，后发生的 $U_z$ 写在左边。几何只能确定最终自旋方向，即物理态的射线；具体整体相位须由下面的矩阵计算确定。

\textbf{2. 显式计算。}记 $c=\cos(\theta/2)$、$s=\sin(\theta/2)$。由 $(-\ii)(-\ii)=-1$、$(-\ii)\ii=1$，
\begin{equation}
  U_y(\theta)=\begin{pmatrix}c&-s\\s&c\end{pmatrix},\qquad
  U_z(\varphi)=\begin{pmatrix}\ee^{-\ii\varphi/2}&0\\0&\ee^{\ii\varphi/2}\end{pmatrix}.
\end{equation}
依次作用得到
\begin{equation}
  U_y(\theta)\ket{z,+}=\begin{pmatrix}c\\s\end{pmatrix},\qquad
  \boxed{U_z(\varphi)U_y(\theta)\ket{z,+}
    =\begin{pmatrix}c\ee^{-\ii\varphi/2}\\s\ee^{\ii\varphi/2}\end{pmatrix}
    =\ket{\hat{\mathbf n},+}.}
\end{equation}
在本题约定下是态矢量的精确相等；与前面的单轴旋转示范相比，不同的旋转操作将此初态送到同一方向，但所得代表矢量可以相差整体相位。
\end{checkedsolution}
